% Self-contained preamble for standalone-rendering of leanification tex files.
% Loaded by `\input{...}` from each subsection's `main.tex`. Concatenated
% from the lecture notes (main.tex preamble + thm.tex + shorts.tex +
% shortsjoris.tex). Regenerated by scaffold/initialize_chapter.py.

%-------Dokumentklasse------------------


%---------packages------------------------
\usepackage[T1]{fontenc}
\usepackage[latin1]{inputenc}
\usepackage[english]{babel}
%\usepackage{amsopn}
\usepackage{stmaryrd}  % for \llbracket, \rrbracket
\SetSymbolFont{stmry}{bold}{U}{stmry}{m}{n}  % to avoid LaTeX Font Warning: Font shape `U/stmry/b/n' undefined (see https://tex.stackexchange.com/questions/106714/stmaryrd-and-boldsymbol-avoid-warnings)
%\newcommand{\changefont}[3]{\fontfamily{#1}\fontseries{#2}\fontshape{#3}\selectfont} 
\usepackage{cite}		%nichtmit natmib zusammen
\usepackage[euler]{textgreek}
\usepackage{enumerate}

\usepackage{amsmath,amssymb,amsfonts}
\usepackage{amsthm}
\usepackage{thm-restate}
%\makeatletter % WORKAROUND FOR BUG WITH RECENT LATEX VERSIONS, SEE https://github.com/muzimuzhi/thmtools/issues/72
%\IfFormatAtLeastTF{2024-11-01}{
%  \renewcommand*\@addtoreset[2]{%
%    \bgroup
%      \edef\aliasctr@@truelist{\aliasctr@follow{#2}}%
%      \let\@elt\relax
%      \expandafter\@cons\aliasctr@@truelist{{#1}}%
%    \egroup
%    \expandafter\xdef\csname theH#1\endcsname{%
%      \expandafter\noexpand\csname theH#2\endcsname.%
%      \noexpand\the\noexpand\value{#1}}%
%  }
%}{}
%\makeatother
\declaretheorem[name=Theorem,numberwithin=section]{thm}
%\usepackage{latexsym,color,mathrsfs,nicefrac,stmaryrd} 
%\usepackage[all]{xy}
%\usepackage{mathabx}  % was used for \asterisk, \Asterisk, but makes too many incompatible changes to other macros, like \not
%\usepackage{mathrsfs}

\usepackage{algorithm}
\usepackage{algpseudocode}
\algnewcommand\algorithmicforeach{\textbf{for each}}
\algdef{S}[FOR]{ForEach}[1]{\algorithmicforeach\ #1\ \algorithmicdo}
% \usepackage{oubraces}  -- not installed on this server

\usepackage{cancel}
\usepackage{centernot}
% \usepackage{bbm}  -- not installed on this server
\usepackage{enumitem}
%\usepackage{wasysym}
\usepackage{tikz}
\usetikzlibrary{arrows,arrows.meta,calc,fit}
\usetikzlibrary{positioning,patterns,matrix}
\usetikzlibrary{decorations.markings,decorations.pathreplacing,decorations.pathmorphing}
\usetikzlibrary{shapes,shapes.arrows,shapes.geometric,shapes.multipart}%,shapes.swigs

\usepackage{graphicx}
\graphicspath{ {./graphics/} }
\usepackage{float}
\usepackage{caption}
\usepackage{subcaption}
\usepackage{verbatim}
\usepackage{bm}
%\usepackage{etoolbox}
%\patchcmd{\figure}{\normalsize}{\normalsize\centering}{}{}
\usepackage{calc}
\usepackage{accents}
% \usepackage{esvect}  -- not installed on this server

\usepackage{listings}
\lstset{language=R,basicstyle=\footnotesize\ttfamily\color{blue!60!red},commentstyle=\itshape,keywordstyle=\bfseries,showstringspaces=false}

\usepackage{hyperref}
\hypersetup{
    colorlinks=true,
%    citecolor=black,
%    filecolor=black,
    linkcolor=blue,
%    urlcolor=black,
    pdfpagemode=UseOutlines,
    pdfencoding=unicode         
}

%-------------conditional annotation marks (Leanify project)------------
\usepackage{etoolbox}
\usepackage{xcolor}
\usepackage{environ}

\newtoggle{showmarks}
\toggletrue{showmarks}
% \togglefalse{showmarks}


\NewEnviron{defmark}{%
  \iftoggle{showmarks}{\textcolor{blue}{Def-start} \BODY\ \textcolor{blue}{Def-end}}{\BODY}%
}

\NewEnviron{claimmark}{%
  \iftoggle{showmarks}{\textcolor{green!50!black}{Claim-start} \BODY\ \textcolor{green!50!black}{Claim-end}}{\BODY}%
}

%-------------theorem envirionments --  abbreviations ------------------

% --- inlined thm.tex ---

% \swapnumbers %Numerierung vor "Satz" anstatt danach%

% \theoremheaderfont{\sf\bfseries}

\theoremstyle{plain}

%\newtheorem*{Satz}{Theorem}
% \newtheorem{sa}{Theorem}[subsubsection]
\newtheorem{sa}{Theorem}[subsection]
%\counterwithin*{sa}{subsection} 
\newtheorem{Thm}[sa]{Theorem}
\newtheorem{DefThm}[sa]{Definition/Theorem}
\newtheorem{Lem}[sa]{Lemma}
\newtheorem{Prp}[sa]{Proposition}
\newtheorem{Cor}[sa]{Corollary}
\newtheorem{Con}[sa]{Conjecture}
\newtheorem{Fct}[sa]{Facts}
\newtheorem{Prn}[sa]{Principle}
\newtheorem{Construction}[sa]{Construction}
\newtheorem{Alg}[sa]{Algorithm}

%\theoremstyle{definition}

\newtheorem{Def}[sa]{Definition}
\newtheorem{DefLem}[sa]{Definition/Lemma}
\newtheorem{NotLem}[sa]{Notation/Lemma}
\newtheorem{Axm}[sa]{Axiom}
\newtheorem{Not}[sa]{Notation}

%\theoremstyle{remark}

\newtheorem{Rem}[sa]{Remark}
\newtheorem{Cau}[sa]{Caution}
\newtheorem{Eg}[sa]{Example}
\newtheorem{Tho}[sa]{Thoughts}
\newtheorem{Exc}[sa]{Exercise}
%\newtheorem{Lit}[sa]{Literature}
\newtheorem{Ques}[sa]{Question}
\newtheorem{Expl}[sa]{Explanation}
\newtheorem{Disc}[sa]{Discussion}
%\newtheorem{Summary}[sa]{Summary}
\newtheorem{Conclusion}[sa]{Conclusion}
\newtheorem{Motivation}[sa]{Motivation}
\newtheorem*{Step*}{Step}
\newtheorem{Step}{Step}
\newtheorem{Ass}[sa]{Assumption}



%\theorembodyfont{\rmfamily}
%\newenvironment{Aufgabe}{\begin{aufg}}{\eoe\end{aufg}}

%% --- inlined shorts.tex ---
\newcommand{\N}{\mathbb{N}}
\newcommand{\Z}{\mathbb{Z}}
\newcommand{\R}{\mathbb{R}}
\newcommand{\Q}{\mathbb{Q}}
\newcommand{\x}{\times}						
\newcommand{\id}{\mathrm{id}}
\newcommand{\Acal}{\mathcal{A}}
\newcommand{\Bcal}{\mathcal{B}}
\newcommand{\Ccal}{\mathcal{C}}
\newcommand{\Dcal}{\mathcal{D}}
\newcommand{\Ecal}{\mathcal{E}}
\newcommand{\Fcal}{\mathcal{F}}
\newcommand{\Gcal}{\mathcal{G}}
\newcommand{\Hcal}{\mathcal{H}}
\newcommand{\Ical}{\mathcal{I}}
\newcommand{\Jcal}{\mathcal{J}}
\newcommand{\Kcal}{\mathcal{K}}
\newcommand{\Lcal}{\mathcal{L}}
\newcommand{\Mcal}{\mathcal{M}}
\newcommand{\Ncal}{\mathcal{N}}
\newcommand{\Ocal}{\mathcal{O}}
\newcommand{\Pcal}{\mathcal{P}}
\newcommand{\Qcal}{\mathcal{Q}}
\newcommand{\Rcal}{\mathcal{R}}
\newcommand{\Scal}{\mathcal{S}}
\newcommand{\Tcal}{\mathcal{T}}
\newcommand{\Ucal}{\mathcal{U}}
\newcommand{\Vcal}{\mathcal{V}}
\newcommand{\Wcal}{\mathcal{W}}
\newcommand{\Xcal}{\mathcal{X}}
\newcommand{\Ycal}{\mathcal{Y}}
\newcommand{\Zcal}{\mathcal{Z}}

\newcommand{\Abf}{\mathbf{A}}
\newcommand{\Bbf}{\mathbf{B}}
\newcommand{\Cbf}{\mathbf{C}}
\newcommand{\Dbf}{\mathbf{D}}
\newcommand{\Ebf}{\mathbf{E}}
\newcommand{\Fbf}{\mathbf{F}}
\newcommand{\Gbf}{\mathbf{G}}
\newcommand{\Hbf}{\mathbf{H}}
\newcommand{\Ibf}{\mathbf{I}}
\newcommand{\Jbf}{\mathbf{J}}
\newcommand{\Kbf}{\mathbf{K}}
\newcommand{\Lbf}{\mathbf{L}}
\newcommand{\Mbf}{\mathbf{M}}
\newcommand{\Nbf}{\mathbf{N}}
\newcommand{\Obf}{\mathbf{O}}
\newcommand{\Pbf}{\mathbf{P}}
\newcommand{\Qbf}{\mathbf{Q}}
\newcommand{\Rbf}{\mathbf{R}}
\newcommand{\Sbf}{\mathbf{S}}
\newcommand{\Tbf}{\mathbf{T}}
\newcommand{\Ubf}{\mathbf{U}}
\newcommand{\Vbf}{\mathbf{V}}
\newcommand{\Wbf}{\mathbf{W}}
\newcommand{\Xbf}{\mathbf{X}}
\newcommand{\Ybf}{\mathbf{Y}}
\newcommand{\Zbf}{\mathbf{Z}}

\newcommand{\tbf}{\mathbf{t}}
\newcommand{\xbf}{\mathbf{x}}

\newcommand{\Qf}{\mathcal{Q}}

\newcommand{\fa}{\mathfrak{a}}
\newcommand{\fb}{\mathfrak{b}}

%\preceq	\succeq  \precsim

\newcommand{\trv}{conditional random variable}
%\newcommand{\trv}{random field}

\newcommand{\track}{track}
\newcommand{\trackable}{trackable}
\newcommand{\trackability}{trackability}
%\newcommand{\Trackable}{Trackable}
\newcommand{\Trackability}{Trackability}

\newcommand{\sm}{\setminus}						
\newcommand{\ins}{\subseteq} 					
\newcommand{\sni}{\supseteq} 					
\newcommand{\cmpl}{\mathsf{c}}

\newcommand{\ol}{\overline}
\newcommand{\ul}{\underline}

\newcommand{\pr}{\mathrm{pr}} 	
\newcommand{\ev}{\mathrm{ev}} 

\newcommand{\srj}{\twoheadrightarrow}
\newcommand{\inj}{\hookrightarrow}
\newcommand{\dshto}{\dashrightarrow}
\newcommand{\dshot}{\dashlefttarrow}

\renewcommand{\Pr}{\mathbb{P}} 	
\newcommand{\Var}{\mathrm{Var}}
\newcommand{\Bias}{\mathrm{Bias}}
\newcommand{\Noise}{\mathrm{Noise}}
\DeclareMathOperator{\Tr}{Tr}

\newcommand{\Msf}{\mathcal{M}}
\newcommand{\Prob}{\mathcal{P}}

	
\newcommand{\E}{\mathbb{E}}
\DeclareMathOperator{\CE}{\mathrm{CE}}
\DeclareMathOperator{\KL}{\mathrm{KL}}
\DeclareMathOperator{\TC}{\mathrm{TC}}
\newcommand{\I}{\mathbbm{1}}
\newcommand{\Pa}{\mathrm{Pa}} 		
\newcommand{\PaD}{\mathrm{PaD}} 		
\newcommand{\Fa}{\mathrm{Fa}} 		
\newcommand{\Pae}{\mathrm{Pae}} 		
\newcommand{\Pai}{\mathrm{Pai}} 		
%\newcommand{\pa}{\mathrm{pa}} 		
\newcommand{\Ch}{\mathrm{Ch}} 	
\newcommand{\Sib}{\mathrm{Sib}} 		
\newcommand{\Adj}{\mathrm{Adj}} 		
\newcommand{\Anc}{\mathrm{Anc}} 		
\newcommand{\Ant}{\mathrm{Ant}}
\newcommand{\Un}{\mathrm{Un}}
\newcommand{\AnCl}{\mathrm{AnCl}} 
\newcommand{\Desc}{\mathrm{Desc}} 	
\newcommand{\Dist}{\mathrm{Dist}} 
\newcommand{\Di}{\mathrm{Di}} 
\newcommand{\Dst}{\mathcal{D}} 
\newcommand{\IDC}{\mathrm{ID}}
\newcommand*{\idc}[1]{\breve{#1}}
\newcommand{\eDist}{\mathrm{eDist}} 
\newcommand{\Pred}{\mathrm{Pred}} 
\newcommand{\Sc}{\mathrm{Sc}}
\newcommand{\MBl}{\mathrm{Mb}}
\newcommand{\NonDesc}{\mathrm{NonDesc}}
\newcommand{\Nb}{\mathrm{Nb}}
%\DeclareMathOperator*{\Indep}{\perp\!\!\!\perp}
\DeclareMathOperator*{\Indep}{\perp\mkern-11mu\perp}
\DeclareMathOperator*{\nIndep}{\not\mkern-2mu\perp\mkern-11mu\perp}
%\DeclareMathOperator*{\nIndep}{\not\Indep}
\DeclareMathOperator*{\Perp}{\perp}
\DeclareMathOperator*{\nPerp}{\not\perp}
\newcommand{\dPerp}{\Perp^d}
\newcommand{\ndPerp}{\nPerp^d}
\newcommand{\mPerp}{\Perp^m}
\newcommand{\nmPerp}{\nPerp^m}
\newcommand{\sPerp}{\Perp^{\sigma}}
\newcommand{\nsPerp}{\nPerp^{\sigma}}
% OLD NOTATION:
%\DeclareMathOperator*{\iPerp}{\perp}
%\DeclareMathOperator*{\niPerp}{\not\perp}
% NEW NOTATION:
%\DeclareMathOperator*{\iPerp}{\mathrel{\tikz \draw [line width=0.5pt] (-0.165em,0.0ex) -- (0.165em,0.0ex) -- (0.165em,0.66em);\hspace{0.33em}}}
%\DeclareMathOperator*{\niPerp}{\not\mkern+2mu\mathrel{\tikz \draw [line width=0.5pt] (0.0em,0.0ex) -- (0.33em,0.0ex) -- (0.33em,0.66em);\hspace{0.33em}}\mkern+2mu}
%\newcommand{\idPerp}{\iPerp^d}
%\newcommand{\nidPerp}{\niPerp^d}
%\newcommand{\isPerp}{\iPerp^{\sigma}}
%\newcommand{\nisPerp}{\niPerp^{\sigma}}
%\newcommand{\imPerp}{\iPerp^m}
%\newcommand{\nimPerp}{\niPerp^m}
% ALTERNATIVE NEW NOTATION
\newcommand{\iPerp}{\Perp^i}
\newcommand{\niPerp}{\nPerp^i}
\newcommand{\idPerp}{\Perp^{id}}
\newcommand{\nidPerp}{\nPerp^{id}}
\newcommand{\isPerp}{\Perp^{i\sigma}}
\newcommand{\nisPerp}{\nPerp^{i\sigma}}
\newcommand{\imPerp}{\Perp^{im}}
\newcommand{\nimPerp}{\nPerp^{im}}
\DeclareMathOperator*{\given}{|}
\DeclareMathOperator*{\kgiven}{\Vert}
\DeclareMathOperator*{\diag}{\mathrm{diag}}
\DeclareMathOperator{\doit}{do}
%\DeclareMathOperator*{\bigdcup}{\dot{\bigcup}}  %error: puts a 9 on top instead of a dot (??)
%\DeclareMathOperator{\dcup}{\dot{\cup}}
%\DeclareMathOperator*{\bigdcup}{\coprod}
%\DeclareMathOperator{\dcup}{\sqcup}
\newcommand{\dcup}{\,\dot{\cup}\,}
\newcommand{\bigdcup}{\mathop{\dot{\bigcup}}}

\DeclareMathOperator*{\argmax}{argmax}
\DeclareMathOperator*{\argmin}{argmin}

\newcommand{\moral}{\mathrm{mor}}	
\newcommand{\marg}{\mathrm{mar}}
%\newcommand{\swig}{\mathrm{o/do}}
\newcommand{\swig}{\mathrm{swig}}
\newcommand{\spl}{\mathrm{split}}
\newcommand{\aug}{\mathrm{aug}}		
\newcommand{\can}{\mathrm{can}}	
\newcommand{\ass}{\mathrm{ass}}		
\newcommand{\augacy}{\mathrm{augacy}}
\newcommand{\acag}{\mathrm{acag}}			
\newcommand{\acy}{\mathrm{acy}}	
\newcommand{\ske}{\mathrm{ske}}	
\newcommand{\cx}{\mathrm{cx}}	
\newcommand{\HEDG}{HEDG}
\newcommand{\Asterisk}{\mathop{\scalebox{1.5}{\raisebox{-0.2ex}{$\ast$}}}}%
\newcommand{\asterisk}{*}

\newcommand{\ReLU}{\mathrm{ReLU}} 
\newcommand{\Lip}{\mathrm{Lip}} 

\newcommand{\lp}{\left ( }
\newcommand{\rp}{\right ) }
\newcommand{\lB}{\left [ }
\newcommand{\rB}{\right ] }
\newcommand{\lC}{\left \{ }
\newcommand{\rC}{\right \} }
\newcommand{\lI}{\left| }
\newcommand{\rI}{\right| }

\newcommand{\st}{\,\middle|\,}

%\newcommand{\tuh}{\rightarrow}
%\newcommand{\hut}{\leftarrow}
%\newcommand{\huh}{\leftrightarrow}
%\newcommand{\tut}{-}
%\newcommand{\tnt}{\diameter}
%\newcommand{\ot}{\leftarrow}
%\newcommand{\oto}{\leftrightarrow}

\newcommand{\arrhead}{{Latex}}
%\newcommand{\arrtail}{{stealth reversed}}  % too visually confusing given al the different edge marks we need
\newcommand{\arrtail}{{}}
\newcommand{\arrstar}{Rays[n=6]}
\newcommand{\arrcirc}{{Circle[open]}}
\newcommand*{\out}[1][]{\mathrel{\tikz [baseline=-0.25ex,\arrcirc-\arrtail, #1] \draw [#1] (0pt,0.5ex) -- (1.3em,0.5ex);}}
\newcommand*{\hut}[1][]{\mathrel{\tikz [baseline=-0.25ex,\arrhead-\arrtail, #1] \draw [#1] (0pt,0.5ex) -- (1.3em,0.5ex);}}
\newcommand*{\tut}[1][]{\mathrel{\tikz [baseline=-0.25ex,\arrtail-\arrtail, #1] \draw [#1] (0pt,0.5ex) -- (1.3em,0.5ex);}}
\newcommand*{\tuh}[1][]{\mathrel{\tikz [baseline=-0.25ex,\arrtail-\arrhead, #1] \draw [#1] (0pt,0.5ex) -- (1.3em,0.5ex);}}
\newcommand*{\tuo}[1][]{\mathrel{\tikz [baseline=-0.25ex,\arrtail-\arrcirc, #1] \draw [#1] (0pt,0.5ex) -- (1.3em,0.5ex);}}
\newcommand*{\huh}[1][]{\mathrel{\tikz [baseline=-0.25ex,\arrhead-\arrhead, #1] \draw [#1] (0pt,0.5ex) -- (1.3em,0.5ex);}}
\newcommand*{\ouo}[1][]{\mathrel{\tikz [baseline=-0.25ex,\arrcirc-\arrcirc, #1] \draw [#1] (0pt,0.5ex) -- (1.3em,0.5ex);}}
\newcommand*{\huo}[1][]{\mathrel{\tikz [baseline=-0.25ex,\arrhead-\arrcirc, #1] \draw [#1] (0pt,0.5ex) -- (1.3em,0.5ex);}}
\newcommand*{\ouh}[1][]{\mathrel{\tikz [baseline=-0.25ex,\arrcirc-\arrhead, #1] \draw [#1] (0pt,0.5ex) -- (1.3em,0.5ex);}}

\newcommand*{\ous}[1][]{\mathrel{\tikz [baseline=-0.25ex,\arrcirc-\arrstar, #1] \draw [#1] (0pt,0.5ex) -- (1.3em,0.5ex);}}
\newcommand*{\hus}[1][]{\mathrel{\tikz [baseline=-0.25ex,\arrhead-\arrstar, #1] \draw [#1] (0pt,0.5ex) -- (1.3em,0.5ex);}}
\newcommand*{\tus}[1][]{\mathrel{\tikz [baseline=-0.25ex,\arrtail-\arrstar, #1] \draw [#1] (0pt,0.5ex) -- (1.3em,0.5ex);}}
\newcommand*{\sut}[1][]{\mathrel{\tikz [baseline=-0.25ex,\arrstar-\arrtail, #1] \draw [#1] (0pt,0.5ex) -- (1.3em,0.5ex);}}
\newcommand*{\suh}[1][]{\mathrel{\tikz [baseline=-0.25ex,\arrstar-\arrhead, #1] \draw [#1] (0pt,0.5ex) -- (1.3em,0.5ex);}}
\newcommand*{\suo}[1][]{\mathrel{\tikz [baseline=-0.25ex,\arrstar-\arrcirc, #1] \draw [#1] (0pt,0.5ex) -- (1.3em,0.5ex);}}
\newcommand*{\sus}[1][]{\mathrel{\tikz [baseline=-0.25ex,\arrstar-\arrstar, #1] \draw [#1] (0pt,0.5ex) -- (1.3em,0.5ex);}}

\newcommand*{\tnt}[1][]{\mathrel{\tikz [baseline=-0.25ex,-, #1] \draw [#1] (0pt,0.5ex) -- node[strike out,draw,-]{} (1.3em,0.5ex);}}
\newcommand*{\toh}[1][]{\mathrel{\tikz [baseline=-0.25ex,-\arrhead, #1] \draw [
decoration={markings, mark=at position 0.4 with {\draw[fill] circle (0.8pt);}},
        postaction={decorate},#1] (0pt,0.5ex) -- (1.3em,0.5ex);}}
\newcommand*{\hot}[1][]{\mathrel{\tikz [baseline=-0.25ex,\arrhead-, #1] \draw [
decoration={markings, mark=at position 0.6 with {\draw[fill] circle (0.8pt);}},
        postaction={decorate},#1] (0pt,0.5ex) -- (1.3em,0.5ex);}}
\newcommand*{\tot}[1][]{\mathrel{\tikz [baseline=-0.25ex,-, #1] \draw [
decoration={markings, mark=at position 0.5 with {\draw[fill] circle (0.8pt);}},
        postaction={decorate},#1] (0pt,0.5ex) -- (1.3em,0.5ex);}}				
\newcommand*{\hoh}[1][]{\mathrel{\tikz [baseline=-0.25ex,\arrhead-\arrhead, #1] \draw [
decoration={markings, mark=at position 0.5 with {\draw[fill] circle (0.8pt);}},
        postaction={decorate},#1] (0pt,0.5ex) -- (1.3em,0.5ex);}}
%\newcommand*{\ouo}[1][]{\mathrel{\tikz [baseline=-0.25ex,-, #1] \draw [decoration={markings, mark=at position 0 with {\draw circle (1pt);}, mark=at position 1 with {\draw circle (1pt);}}, postaction={decorate},#1] (0pt,0.5ex) -- (1.3em,0.5ex);}}

%further arrow tips: Arc Barb, Bar, Bracket, Hooks, Parenthesis, Tee Bar, Circle, Diamond, Rectangle, Square, Rays[n=6]

%\newcommand*{\rars}{\begin{array}{c}\huh\\[-10pt]\tuh\end{array}}
%\newcommand*{\lars}{\begin{array}{c}\huh\\[-10pt]\hut\end{array}}
%\newcommand*{\allars}{\begin{array}{c}\huh\\[-10pt]\tuh\\[-10pt]\hut\\[-10pt]\tut\end{array}}
%\newcommand*{\dars}{\begin{array}{c}\huh\\[-10pt]\tuh\\[-10pt]\hut\end{array}}


\newcommand*{\matb}[1]{\begin{bmatrix}#1\end{bmatrix}}
\newcommand*{\mat}[1]{\begin{pmatrix}#1\end{pmatrix}}


\makeatletter
\newcommand{\xRightarrow}[2][]{\ext@arrow 0359\Rightarrowfill@{#1}{#2}}
\makeatother


\newcommand{\PF}[1]{{\color{orange}{PF: #1}}}
%\newcommand{\PF}[1]{}


%\tikzstyle{ndint} = [draw, line width=1pt, color=teal, shape=rectangle, minimum size=20pt,inner sep=0pt, text=black]
%\tikzstyle{ndout} = [draw, line width=1pt, shape=circle, color=blue, minimum size=20pt,inner sep=0pt, text=black]
%\tikzstyle{ndlat} = [draw, line width=1pt, shape=circle, color=red, minimum size=20pt,inner sep=0pt, text=black]
\tikzstyle{ndint} = [draw, semithick, shape=rectangle, minimum size=20pt,inner sep=0pt]
\tikzstyle{ndout} = [draw, semithick, shape=circle, minimum size=20pt,inner sep=0pt]
\tikzstyle{ndlat} = [draw, semithick, shape=circle, minimum size=20pt,inner sep=0pt, fill=lightgray]
%\tikzstyle{ndvst} = [draw, ultra thick, shape=swig vsplit, swig vsplit={gap=5pt,line color right=teal, line color left=blue}, inner sep=1.2pt, minimum size=7mm]
%\tikzstyle{ndhst} = [draw, ultra thick, shape=swig hsplit, swig hsplit={gap=5pt,line color lower=teal, line color upper=blue}, inner sep=1.2pt, minimum size=7mm]
\tikzstyle{arint} = [style={->,>=Latex,thick}] %teal
\tikzstyle{arout} = [style={->,>=Latex,thick}] %blue
\tikzstyle{arlout} = [style={->,>=Latex,thick}] %red
\tikzstyle{arlat} = [style={<->,>=Latex,thick}] %red

\tikzstyle{out} = [style={o->,>=stealth',thick,draw,fill}] %this gives a tikz notation clash, for the out= parameter to specify the outgoing angle of an edge :(
\tikzstyle{hut} = [style={<-,>=Latex,style=semithick}]
\tikzstyle{tut} = [style=semithick]
\tikzstyle{tuh} = [style={->,>=Latex,style=semithick}]
\tikzstyle{tuo} = [style={-o,style=semithick}]
\tikzstyle{huh} = [style={<->,>=Latex,style=semithick}]
\tikzstyle{ouo} = [style={o-o,style=semithick}]
\tikzstyle{huo} = [style={<-o,>=Latex,style=semithick}]
\tikzstyle{ouh} = [style={o->,>=Latex,style=semithick}]
\tikzstyle{ous} = [style={o-{Rays[n=6]},style=semithick}]
\tikzstyle{hus} = [style={<-{Rays[n=6]},>=Latex,style=semithick}]
\tikzstyle{tus} = [style={-{Rays[n=6]},style=semithick}]
\tikzstyle{sut} = [style={{Rays[n=6]}-,style=semithick}]
\tikzstyle{suh} = [style={{Rays[n=6]}->,>=Latex,style=semithick}]
\tikzstyle{suo} = [style={{Rays[n=6]}-o,style=semithick}]
\tikzstyle{sus} = [style={{Rays[n=6]}-{Rays[n=6]},style=semithick}]

% --- end shorts.tex ---

% --- end thm.tex ---
% --- inlined shorts.tex ---
\newcommand{\N}{\mathbb{N}}
\newcommand{\Z}{\mathbb{Z}}
\newcommand{\R}{\mathbb{R}}
\newcommand{\Q}{\mathbb{Q}}
\newcommand{\x}{\times}						
\newcommand{\id}{\mathrm{id}}
\newcommand{\Acal}{\mathcal{A}}
\newcommand{\Bcal}{\mathcal{B}}
\newcommand{\Ccal}{\mathcal{C}}
\newcommand{\Dcal}{\mathcal{D}}
\newcommand{\Ecal}{\mathcal{E}}
\newcommand{\Fcal}{\mathcal{F}}
\newcommand{\Gcal}{\mathcal{G}}
\newcommand{\Hcal}{\mathcal{H}}
\newcommand{\Ical}{\mathcal{I}}
\newcommand{\Jcal}{\mathcal{J}}
\newcommand{\Kcal}{\mathcal{K}}
\newcommand{\Lcal}{\mathcal{L}}
\newcommand{\Mcal}{\mathcal{M}}
\newcommand{\Ncal}{\mathcal{N}}
\newcommand{\Ocal}{\mathcal{O}}
\newcommand{\Pcal}{\mathcal{P}}
\newcommand{\Qcal}{\mathcal{Q}}
\newcommand{\Rcal}{\mathcal{R}}
\newcommand{\Scal}{\mathcal{S}}
\newcommand{\Tcal}{\mathcal{T}}
\newcommand{\Ucal}{\mathcal{U}}
\newcommand{\Vcal}{\mathcal{V}}
\newcommand{\Wcal}{\mathcal{W}}
\newcommand{\Xcal}{\mathcal{X}}
\newcommand{\Ycal}{\mathcal{Y}}
\newcommand{\Zcal}{\mathcal{Z}}

\newcommand{\Abf}{\mathbf{A}}
\newcommand{\Bbf}{\mathbf{B}}
\newcommand{\Cbf}{\mathbf{C}}
\newcommand{\Dbf}{\mathbf{D}}
\newcommand{\Ebf}{\mathbf{E}}
\newcommand{\Fbf}{\mathbf{F}}
\newcommand{\Gbf}{\mathbf{G}}
\newcommand{\Hbf}{\mathbf{H}}
\newcommand{\Ibf}{\mathbf{I}}
\newcommand{\Jbf}{\mathbf{J}}
\newcommand{\Kbf}{\mathbf{K}}
\newcommand{\Lbf}{\mathbf{L}}
\newcommand{\Mbf}{\mathbf{M}}
\newcommand{\Nbf}{\mathbf{N}}
\newcommand{\Obf}{\mathbf{O}}
\newcommand{\Pbf}{\mathbf{P}}
\newcommand{\Qbf}{\mathbf{Q}}
\newcommand{\Rbf}{\mathbf{R}}
\newcommand{\Sbf}{\mathbf{S}}
\newcommand{\Tbf}{\mathbf{T}}
\newcommand{\Ubf}{\mathbf{U}}
\newcommand{\Vbf}{\mathbf{V}}
\newcommand{\Wbf}{\mathbf{W}}
\newcommand{\Xbf}{\mathbf{X}}
\newcommand{\Ybf}{\mathbf{Y}}
\newcommand{\Zbf}{\mathbf{Z}}

\newcommand{\tbf}{\mathbf{t}}
\newcommand{\xbf}{\mathbf{x}}

\newcommand{\Qf}{\mathcal{Q}}

\newcommand{\fa}{\mathfrak{a}}
\newcommand{\fb}{\mathfrak{b}}

%\preceq	\succeq  \precsim

\newcommand{\trv}{conditional random variable}
%\newcommand{\trv}{random field}

\newcommand{\track}{track}
\newcommand{\trackable}{trackable}
\newcommand{\trackability}{trackability}
%\newcommand{\Trackable}{Trackable}
\newcommand{\Trackability}{Trackability}

\newcommand{\sm}{\setminus}						
\newcommand{\ins}{\subseteq} 					
\newcommand{\sni}{\supseteq} 					
\newcommand{\cmpl}{\mathsf{c}}

\newcommand{\ol}{\overline}
\newcommand{\ul}{\underline}

\newcommand{\pr}{\mathrm{pr}} 	
\newcommand{\ev}{\mathrm{ev}} 

\newcommand{\srj}{\twoheadrightarrow}
\newcommand{\inj}{\hookrightarrow}
\newcommand{\dshto}{\dashrightarrow}
\newcommand{\dshot}{\dashlefttarrow}

\renewcommand{\Pr}{\mathbb{P}} 	
\newcommand{\Var}{\mathrm{Var}}
\newcommand{\Bias}{\mathrm{Bias}}
\newcommand{\Noise}{\mathrm{Noise}}
\DeclareMathOperator{\Tr}{Tr}

\newcommand{\Msf}{\mathcal{M}}
\newcommand{\Prob}{\mathcal{P}}

	
\newcommand{\E}{\mathbb{E}}
\DeclareMathOperator{\CE}{\mathrm{CE}}
\DeclareMathOperator{\KL}{\mathrm{KL}}
\DeclareMathOperator{\TC}{\mathrm{TC}}
\newcommand{\I}{\mathbbm{1}}
\newcommand{\Pa}{\mathrm{Pa}} 		
\newcommand{\PaD}{\mathrm{PaD}} 		
\newcommand{\Fa}{\mathrm{Fa}} 		
\newcommand{\Pae}{\mathrm{Pae}} 		
\newcommand{\Pai}{\mathrm{Pai}} 		
%\newcommand{\pa}{\mathrm{pa}} 		
\newcommand{\Ch}{\mathrm{Ch}} 	
\newcommand{\Sib}{\mathrm{Sib}} 		
\newcommand{\Adj}{\mathrm{Adj}} 		
\newcommand{\Anc}{\mathrm{Anc}} 		
\newcommand{\Ant}{\mathrm{Ant}}
\newcommand{\Un}{\mathrm{Un}}
\newcommand{\AnCl}{\mathrm{AnCl}} 
\newcommand{\Desc}{\mathrm{Desc}} 	
\newcommand{\Dist}{\mathrm{Dist}} 
\newcommand{\Di}{\mathrm{Di}} 
\newcommand{\Dst}{\mathcal{D}} 
\newcommand{\IDC}{\mathrm{ID}}
\newcommand*{\idc}[1]{\breve{#1}}
\newcommand{\eDist}{\mathrm{eDist}} 
\newcommand{\Pred}{\mathrm{Pred}} 
\newcommand{\Sc}{\mathrm{Sc}}
\newcommand{\MBl}{\mathrm{Mb}}
\newcommand{\NonDesc}{\mathrm{NonDesc}}
\newcommand{\Nb}{\mathrm{Nb}}
%\DeclareMathOperator*{\Indep}{\perp\!\!\!\perp}
\DeclareMathOperator*{\Indep}{\perp\mkern-11mu\perp}
\DeclareMathOperator*{\nIndep}{\not\mkern-2mu\perp\mkern-11mu\perp}
%\DeclareMathOperator*{\nIndep}{\not\Indep}
\DeclareMathOperator*{\Perp}{\perp}
\DeclareMathOperator*{\nPerp}{\not\perp}
\newcommand{\dPerp}{\Perp^d}
\newcommand{\ndPerp}{\nPerp^d}
\newcommand{\mPerp}{\Perp^m}
\newcommand{\nmPerp}{\nPerp^m}
\newcommand{\sPerp}{\Perp^{\sigma}}
\newcommand{\nsPerp}{\nPerp^{\sigma}}
% OLD NOTATION:
%\DeclareMathOperator*{\iPerp}{\perp}
%\DeclareMathOperator*{\niPerp}{\not\perp}
% NEW NOTATION:
%\DeclareMathOperator*{\iPerp}{\mathrel{\tikz \draw [line width=0.5pt] (-0.165em,0.0ex) -- (0.165em,0.0ex) -- (0.165em,0.66em);\hspace{0.33em}}}
%\DeclareMathOperator*{\niPerp}{\not\mkern+2mu\mathrel{\tikz \draw [line width=0.5pt] (0.0em,0.0ex) -- (0.33em,0.0ex) -- (0.33em,0.66em);\hspace{0.33em}}\mkern+2mu}
%\newcommand{\idPerp}{\iPerp^d}
%\newcommand{\nidPerp}{\niPerp^d}
%\newcommand{\isPerp}{\iPerp^{\sigma}}
%\newcommand{\nisPerp}{\niPerp^{\sigma}}
%\newcommand{\imPerp}{\iPerp^m}
%\newcommand{\nimPerp}{\niPerp^m}
% ALTERNATIVE NEW NOTATION
\newcommand{\iPerp}{\Perp^i}
\newcommand{\niPerp}{\nPerp^i}
\newcommand{\idPerp}{\Perp^{id}}
\newcommand{\nidPerp}{\nPerp^{id}}
\newcommand{\isPerp}{\Perp^{i\sigma}}
\newcommand{\nisPerp}{\nPerp^{i\sigma}}
\newcommand{\imPerp}{\Perp^{im}}
\newcommand{\nimPerp}{\nPerp^{im}}
\DeclareMathOperator*{\given}{|}
\DeclareMathOperator*{\kgiven}{\Vert}
\DeclareMathOperator*{\diag}{\mathrm{diag}}
\DeclareMathOperator{\doit}{do}
%\DeclareMathOperator*{\bigdcup}{\dot{\bigcup}}  %error: puts a 9 on top instead of a dot (??)
%\DeclareMathOperator{\dcup}{\dot{\cup}}
%\DeclareMathOperator*{\bigdcup}{\coprod}
%\DeclareMathOperator{\dcup}{\sqcup}
\newcommand{\dcup}{\,\dot{\cup}\,}
\newcommand{\bigdcup}{\mathop{\dot{\bigcup}}}

\DeclareMathOperator*{\argmax}{argmax}
\DeclareMathOperator*{\argmin}{argmin}

\newcommand{\moral}{\mathrm{mor}}	
\newcommand{\marg}{\mathrm{mar}}
%\newcommand{\swig}{\mathrm{o/do}}
\newcommand{\swig}{\mathrm{swig}}
\newcommand{\spl}{\mathrm{split}}
\newcommand{\aug}{\mathrm{aug}}		
\newcommand{\can}{\mathrm{can}}	
\newcommand{\ass}{\mathrm{ass}}		
\newcommand{\augacy}{\mathrm{augacy}}
\newcommand{\acag}{\mathrm{acag}}			
\newcommand{\acy}{\mathrm{acy}}	
\newcommand{\ske}{\mathrm{ske}}	
\newcommand{\cx}{\mathrm{cx}}	
\newcommand{\HEDG}{HEDG}
\newcommand{\Asterisk}{\mathop{\scalebox{1.5}{\raisebox{-0.2ex}{$\ast$}}}}%
\newcommand{\asterisk}{*}

\newcommand{\ReLU}{\mathrm{ReLU}} 
\newcommand{\Lip}{\mathrm{Lip}} 

\newcommand{\lp}{\left ( }
\newcommand{\rp}{\right ) }
\newcommand{\lB}{\left [ }
\newcommand{\rB}{\right ] }
\newcommand{\lC}{\left \{ }
\newcommand{\rC}{\right \} }
\newcommand{\lI}{\left| }
\newcommand{\rI}{\right| }

\newcommand{\st}{\,\middle|\,}

%\newcommand{\tuh}{\rightarrow}
%\newcommand{\hut}{\leftarrow}
%\newcommand{\huh}{\leftrightarrow}
%\newcommand{\tut}{-}
%\newcommand{\tnt}{\diameter}
%\newcommand{\ot}{\leftarrow}
%\newcommand{\oto}{\leftrightarrow}

\newcommand{\arrhead}{{Latex}}
%\newcommand{\arrtail}{{stealth reversed}}  % too visually confusing given al the different edge marks we need
\newcommand{\arrtail}{{}}
\newcommand{\arrstar}{Rays[n=6]}
\newcommand{\arrcirc}{{Circle[open]}}
\newcommand*{\out}[1][]{\mathrel{\tikz [baseline=-0.25ex,\arrcirc-\arrtail, #1] \draw [#1] (0pt,0.5ex) -- (1.3em,0.5ex);}}
\newcommand*{\hut}[1][]{\mathrel{\tikz [baseline=-0.25ex,\arrhead-\arrtail, #1] \draw [#1] (0pt,0.5ex) -- (1.3em,0.5ex);}}
\newcommand*{\tut}[1][]{\mathrel{\tikz [baseline=-0.25ex,\arrtail-\arrtail, #1] \draw [#1] (0pt,0.5ex) -- (1.3em,0.5ex);}}
\newcommand*{\tuh}[1][]{\mathrel{\tikz [baseline=-0.25ex,\arrtail-\arrhead, #1] \draw [#1] (0pt,0.5ex) -- (1.3em,0.5ex);}}
\newcommand*{\tuo}[1][]{\mathrel{\tikz [baseline=-0.25ex,\arrtail-\arrcirc, #1] \draw [#1] (0pt,0.5ex) -- (1.3em,0.5ex);}}
\newcommand*{\huh}[1][]{\mathrel{\tikz [baseline=-0.25ex,\arrhead-\arrhead, #1] \draw [#1] (0pt,0.5ex) -- (1.3em,0.5ex);}}
\newcommand*{\ouo}[1][]{\mathrel{\tikz [baseline=-0.25ex,\arrcirc-\arrcirc, #1] \draw [#1] (0pt,0.5ex) -- (1.3em,0.5ex);}}
\newcommand*{\huo}[1][]{\mathrel{\tikz [baseline=-0.25ex,\arrhead-\arrcirc, #1] \draw [#1] (0pt,0.5ex) -- (1.3em,0.5ex);}}
\newcommand*{\ouh}[1][]{\mathrel{\tikz [baseline=-0.25ex,\arrcirc-\arrhead, #1] \draw [#1] (0pt,0.5ex) -- (1.3em,0.5ex);}}

\newcommand*{\ous}[1][]{\mathrel{\tikz [baseline=-0.25ex,\arrcirc-\arrstar, #1] \draw [#1] (0pt,0.5ex) -- (1.3em,0.5ex);}}
\newcommand*{\hus}[1][]{\mathrel{\tikz [baseline=-0.25ex,\arrhead-\arrstar, #1] \draw [#1] (0pt,0.5ex) -- (1.3em,0.5ex);}}
\newcommand*{\tus}[1][]{\mathrel{\tikz [baseline=-0.25ex,\arrtail-\arrstar, #1] \draw [#1] (0pt,0.5ex) -- (1.3em,0.5ex);}}
\newcommand*{\sut}[1][]{\mathrel{\tikz [baseline=-0.25ex,\arrstar-\arrtail, #1] \draw [#1] (0pt,0.5ex) -- (1.3em,0.5ex);}}
\newcommand*{\suh}[1][]{\mathrel{\tikz [baseline=-0.25ex,\arrstar-\arrhead, #1] \draw [#1] (0pt,0.5ex) -- (1.3em,0.5ex);}}
\newcommand*{\suo}[1][]{\mathrel{\tikz [baseline=-0.25ex,\arrstar-\arrcirc, #1] \draw [#1] (0pt,0.5ex) -- (1.3em,0.5ex);}}
\newcommand*{\sus}[1][]{\mathrel{\tikz [baseline=-0.25ex,\arrstar-\arrstar, #1] \draw [#1] (0pt,0.5ex) -- (1.3em,0.5ex);}}

\newcommand*{\tnt}[1][]{\mathrel{\tikz [baseline=-0.25ex,-, #1] \draw [#1] (0pt,0.5ex) -- node[strike out,draw,-]{} (1.3em,0.5ex);}}
\newcommand*{\toh}[1][]{\mathrel{\tikz [baseline=-0.25ex,-\arrhead, #1] \draw [
decoration={markings, mark=at position 0.4 with {\draw[fill] circle (0.8pt);}},
        postaction={decorate},#1] (0pt,0.5ex) -- (1.3em,0.5ex);}}
\newcommand*{\hot}[1][]{\mathrel{\tikz [baseline=-0.25ex,\arrhead-, #1] \draw [
decoration={markings, mark=at position 0.6 with {\draw[fill] circle (0.8pt);}},
        postaction={decorate},#1] (0pt,0.5ex) -- (1.3em,0.5ex);}}
\newcommand*{\tot}[1][]{\mathrel{\tikz [baseline=-0.25ex,-, #1] \draw [
decoration={markings, mark=at position 0.5 with {\draw[fill] circle (0.8pt);}},
        postaction={decorate},#1] (0pt,0.5ex) -- (1.3em,0.5ex);}}				
\newcommand*{\hoh}[1][]{\mathrel{\tikz [baseline=-0.25ex,\arrhead-\arrhead, #1] \draw [
decoration={markings, mark=at position 0.5 with {\draw[fill] circle (0.8pt);}},
        postaction={decorate},#1] (0pt,0.5ex) -- (1.3em,0.5ex);}}
%\newcommand*{\ouo}[1][]{\mathrel{\tikz [baseline=-0.25ex,-, #1] \draw [decoration={markings, mark=at position 0 with {\draw circle (1pt);}, mark=at position 1 with {\draw circle (1pt);}}, postaction={decorate},#1] (0pt,0.5ex) -- (1.3em,0.5ex);}}

%further arrow tips: Arc Barb, Bar, Bracket, Hooks, Parenthesis, Tee Bar, Circle, Diamond, Rectangle, Square, Rays[n=6]

%\newcommand*{\rars}{\begin{array}{c}\huh\\[-10pt]\tuh\end{array}}
%\newcommand*{\lars}{\begin{array}{c}\huh\\[-10pt]\hut\end{array}}
%\newcommand*{\allars}{\begin{array}{c}\huh\\[-10pt]\tuh\\[-10pt]\hut\\[-10pt]\tut\end{array}}
%\newcommand*{\dars}{\begin{array}{c}\huh\\[-10pt]\tuh\\[-10pt]\hut\end{array}}


\newcommand*{\matb}[1]{\begin{bmatrix}#1\end{bmatrix}}
\newcommand*{\mat}[1]{\begin{pmatrix}#1\end{pmatrix}}


\makeatletter
\newcommand{\xRightarrow}[2][]{\ext@arrow 0359\Rightarrowfill@{#1}{#2}}
\makeatother


\newcommand{\PF}[1]{{\color{orange}{PF: #1}}}
%\newcommand{\PF}[1]{}


%\tikzstyle{ndint} = [draw, line width=1pt, color=teal, shape=rectangle, minimum size=20pt,inner sep=0pt, text=black]
%\tikzstyle{ndout} = [draw, line width=1pt, shape=circle, color=blue, minimum size=20pt,inner sep=0pt, text=black]
%\tikzstyle{ndlat} = [draw, line width=1pt, shape=circle, color=red, minimum size=20pt,inner sep=0pt, text=black]
\tikzstyle{ndint} = [draw, semithick, shape=rectangle, minimum size=20pt,inner sep=0pt]
\tikzstyle{ndout} = [draw, semithick, shape=circle, minimum size=20pt,inner sep=0pt]
\tikzstyle{ndlat} = [draw, semithick, shape=circle, minimum size=20pt,inner sep=0pt, fill=lightgray]
%\tikzstyle{ndvst} = [draw, ultra thick, shape=swig vsplit, swig vsplit={gap=5pt,line color right=teal, line color left=blue}, inner sep=1.2pt, minimum size=7mm]
%\tikzstyle{ndhst} = [draw, ultra thick, shape=swig hsplit, swig hsplit={gap=5pt,line color lower=teal, line color upper=blue}, inner sep=1.2pt, minimum size=7mm]
\tikzstyle{arint} = [style={->,>=Latex,thick}] %teal
\tikzstyle{arout} = [style={->,>=Latex,thick}] %blue
\tikzstyle{arlout} = [style={->,>=Latex,thick}] %red
\tikzstyle{arlat} = [style={<->,>=Latex,thick}] %red

\tikzstyle{out} = [style={o->,>=stealth',thick,draw,fill}] %this gives a tikz notation clash, for the out= parameter to specify the outgoing angle of an edge :(
\tikzstyle{hut} = [style={<-,>=Latex,style=semithick}]
\tikzstyle{tut} = [style=semithick]
\tikzstyle{tuh} = [style={->,>=Latex,style=semithick}]
\tikzstyle{tuo} = [style={-o,style=semithick}]
\tikzstyle{huh} = [style={<->,>=Latex,style=semithick}]
\tikzstyle{ouo} = [style={o-o,style=semithick}]
\tikzstyle{huo} = [style={<-o,>=Latex,style=semithick}]
\tikzstyle{ouh} = [style={o->,>=Latex,style=semithick}]
\tikzstyle{ous} = [style={o-{Rays[n=6]},style=semithick}]
\tikzstyle{hus} = [style={<-{Rays[n=6]},>=Latex,style=semithick}]
\tikzstyle{tus} = [style={-{Rays[n=6]},style=semithick}]
\tikzstyle{sut} = [style={{Rays[n=6]}-,style=semithick}]
\tikzstyle{suh} = [style={{Rays[n=6]}->,>=Latex,style=semithick}]
\tikzstyle{suo} = [style={{Rays[n=6]}-o,style=semithick}]
\tikzstyle{sus} = [style={{Rays[n=6]}-{Rays[n=6]},style=semithick}]

% --- end shorts.tex ---
% --- inlined shortsjoris.tex ---
\newcommand\RN{\R}
\newcommand\B[1]{\bm{#1}}
%\renewcommand\C[1]{\mathcal{#1}}     %PF: suddenly got many errors that \C was already defined ???
%\newcommand\C[1]{\mathcal{#1}}     %PF: suddenly got many errors that \C was NOT defined ???
% in some newer latex distributions, \C is already defined, in others, it is not...! we override it in case it is already defined
\makeatletter
\@ifundefined{C}
{% \C not defined
\newcommand\C[1]{\mathcal{#1}}     
}
{% \C defined
\renewcommand\C[1]{\mathcal{#1}}
}
\makeatother

\newcommand\BC[1]{\bm{\mathcal{#1}}}
\renewcommand\S[1]{\mathscr{#1}}
\newcommand\eref[1]{(\ref{#1})}
\newcommand\Prb{\mathbb{P}}
\newcommand\cond{\mid}
\newcommand\compose{\mathrm{comp}}
%\newcommand\Joris[1]{{\color{purple}JM: #1}}
%\newcommand\JorisTODO[1]{{\color{red}JM TODO: #1}}
\newcommand\JorisTODO[1]{}
\newcommand\Joris[1]{}
%\newcommand\JorisThoughts[1]{\color{orange}JM: #1}
\newcommand\JorisThoughts[1]{}
\newcommand\Claude[1]{{\color{blue}Claude: #1}}
%\newcommand\Claude[1]{}
\newenvironment{topalignedcolumn}[1]{\begin{minipage}[t]{#1}\vspace{0pt}\relax}{\end{minipage}}
\newcommand\define[1]{\textbf{#1}}
\newcommand\ind{\I}
\newcommand\foralmostall{\hspace*{0.1em}{\scalebox{1}[1.156]{$\vee$}}\hspace*{-0.74em}\raisebox{0.5ex}{\scalebox{1}[0.8]{$\sim$}}}
\newcommand\notimplies{\centernot\implies}
\newcommand\as{\text{-a.s.}}
\newcommand\mas{\text{-m.a.s.}}
\newcommand\TODOmas{\Joris{Use m.a.s.? }}

\newcommand\twin{\mathrm{twin}}
\newcommand\skel{\mathrm{skel}}
\newcommand{\ot}{\mathrel{\leftarrow}}
%\newcommand{\oto}{\mathrel{\leftrightarrow}}
%\newcommand{\ots}{\mathrel{{\leftarrow\mkern-11mu\ast}}}
%\newcommand{\otc}{\mathrel{{\leftarrow\mkern-8mu\circ}}}
%\newcommand{\sto}{\mathrel{{\ast\mkern-11mu\to}}}
%\newcommand{\stt}{\mathrel{{\ast\mkern-11mu\relbar\mkern-9mu\relbar}}}
%\newcommand{\tts}{\mathrel{{\relbar\mkern-11mu\relbar\mkern-11mu\ast}}}
%\newcommand{\sts}{\mathrel{{\ast\mkern-11mu\relbar\mkern-9mu\relbar\mkern-11mu\ast}}}
%\newcommand{\ctc}{\mathrel{{\circ\mkern-8mu\relbar\mkern-8mu\circ}}}
%\newcommand{\cto}{\mathrel{{\circ\mkern-8mu\rightarrow}}}
%\newcommand{\cts}{\mathrel{{\circ\mkern-8mu\relbar\mkern-9mu\relbar\mkern-11mu\ast}}}
%\newcommand{\stc}{\mathrel{{\ast\mkern-10mu\relbar\mkern-9mu\relbar\mkern-8mu\circ}}}
%\newcommand{\ctt}{\mathrel{{\circ\mkern-8mu\relbar\mkern-9mu\relbar}}}
%\newcommand{\ttt}{\mathrel{{\relbar\mkern-9mu\relbar}}}
%\newcommand{\ttc}{\mathrel{{\relbar\mkern-9mu\relbar\mkern-8mu\circ}}}
%\newcommand\otto{\mathrel{{\Relbar\mkern-9mu\Relbar}}}

\newcommand\Sin{J}
\newcommand\Send{V}
\newcommand\Srnd{W}
\newcommand\ST{T}
\newcommand\xJ{x_{\Sin}}
\newcommand\xV{x_{\Send}}
\newcommand\xW{x_{\Srnd}}
\newcommand\xT{x_{\ST}}
\newcommand\XJ{X_{\Sin}}
\newcommand\XV{X_{\Send}}
\newcommand\XW{X_{\Srnd}}
\newcommand\XT{X_{\ST}}
\newcommand\CX{\C{X}}
\newcommand\CXJ{\CX_{\Sin}}
\newcommand\CXV{\CX_{\Send}}
\newcommand\CXW{\CX_{\Srnd}}
\newcommand\CXT{\CX_{\ST}}

\newcommand{\Exp}{\mathbb{E}}
\newcommand{\Cov}{\mathbb{C}\mathrm{ov}}
\newcommand{\Vaf}{\mathbb{V}\mathrm{ar}}
\newcommand{\dbtilde}[1]{\accentset{\approx}{#1}}

\newcommand\IM{\mathrm{IM}}
\newcommand{\MAG}{\mathrm{MAG}}
\newcommand{\PAG}{\mathrm{PAG}}
\newcommand{\DMAG}{\mathrm{DMAG}}
\newcommand{\DPAG}{\mathrm{DPAG}}
\newcommand{\CDPAG}{\mathrm{CDPAG}}
\newcommand{\sigmaPAG}{{\mathrm{PAG}^\sigma}}
\newcommand{\sigmaMAG}{{\mathrm{MAG}^\sigma}}
\newcommand{\sepset}{\mathtt{sepset}}
\newcommand{\separable}{\mathtt{separable}}
\newcommand{\FCI}{\mathtt{FCI}}
\newcommand{\SEP}{\mathtt{SEP}}
\newcommand{\PosSEP}{\mathtt{posSEP}}

\newcommand{\repar}[2]{\mathrm{repar}({#1},{#2})}

\newcommand{\lt}{\lp} % or: \langle
\newcommand{\rt}{\rp} % or: \rangle
\newcommand{\ltuple}{(}
\newcommand{\rtuple}{)}

\newcommand{\Break}{\State \textbf{break}}
\newcommand{\Input}{\State \textbf{input }}
\newcommand{\Output}{\State \textbf{output }}
\newcommand{\Sample}{\State \textbf{sample }}

\newcommand{\observational}{observable }

% --- end shortsjoris.tex ---
\renewcommand\eminnershape{\itshape\bfseries}

%--------------Renumbering of sections, chapters, subsections----------------
%\usepackage{chngcntr}

%\renewcommand{\thesa}{\arabic{section}.\arabic{subsection}.\arabic{sa}}
%%\renewcommand{\thesa}{\arabic{subsection}.\arabic{subsubsection}.\arabic{sa}}
%%\renewcommand{\thechapter}{\arabic{chapter}}
%\renewcommand{\thesection}{\arabic{section}}
%\renewcommand{\thesubsection}{\arabic{subsection}}
%\renewcommand{\thesubsection}{\arabic{section}.\arabic{subsection}}
%%\counterwithout{section}{chapter} %make sections independent of chapters
%%\counterwithout{subsection}{section} %make sections independent of chapters
%%\counterwithout{section}{subsection}

% --------------------------------------------------
% Throughout the text \subsection* is used which resets the \subsubsection counts.
% The following fixes this issue:
% --------------------------------------------------
% \counterwithout{subsubsection}{subsection}
% \newcounter{customsubsub}
% \renewcommand{\thesa}{\arabic{subsection}.\arabic{customsubsub}.\arabic{sa}}
% \renewcommand{\thesubsubsection}{\arabic{subsection}.\arabic{customsubsub}}
% \let\origsubsection\subsection
% \let\origsubsubsection\subsubsection
% \makeatletter
% \renewcommand{\subsection}{\@ifstar{\origsubsection*}{\setcounter{customsubsub}{0}\origsubsection}}
% \renewcommand{\subsubsection}{\stepcounter{customsubsub}\origsubsubsection}
% \makeatother


%--------------------------------------------------
\setcounter{tocdepth}{5}
\setcounter{secnumdepth}{3}

%-------------------
